% Transcription — fonds Alexandre Grothendieck, shelfmark 63, batch 3 (pages 41-60).
% Established from the facsimile archives/batches/63/batch-03.pdf (a copy of
% 63.pdf fetched through the project relay, no cover sheet: PDF page k =
% archivists' page 40+k), per the skill transcribe-grothendieck. The folder is
% a hundred pages in five batches; this is the third. Aligned with the
% decisions of batches 1 and 2 (transcripts/63/batch-01.fr.tex,
% batch-02.fr.tex). Folder 62 (transcripts/62/) is the notational reference
% for the group « Courbes elliptiques » (62-64).
%
% Pass: Opus 5.5 (claude-opus-5-5), 2026-09-26 — first pass, unchecked against the pages by a human.
%
% Pages skipped: 55, a blank cream sheet.
%
% Nothing in this batch is transcribed in full: no page carries mathematics
% in his hand. Pages given one \note each, not re-set (someone else's text):
% - Pages 41-54: typed pages « atkin 2 » ... « atkin 15 » of the typescript
%   copy of A. O. L. Atkin, « Congruences for modular forms », whose first
%   page (« atkin 1 », « (1968) » in ink) is page 40 of batch 2; the text
%   runs on without a break from there (p. 41 opens « with Dr.
%   R.F.Churchhouse ... »), through §2 (Ramanujan congruences for q <= 11),
%   §3 (q = 13, Newman, the M- and R-results), §4 (q = 17, 19, 23), §5
%   (q >= 29, GF(29^2)), §6 (analogy with Hecke theory), §7 (non-Ramanujan
%   congruences for q <= 11) and §8 (notes on computation), ending on typed
%   p. 15 (p. 54) with no bibliography in the copy. Evidence that it is not
%   his: the author's name on the typed header of every leaf; first-person
%   narrative that is Atkin's (« After my appointment to Atlas in 1964 »,
%   « In 1962 O'Brien and I »); English throughout; the formulas are filled
%   in by hand in a rounded, upright English hand nothing like his.
% - Pages 56-60: the first five leaves of a second typescript, A. O. L.
%   Atkin, « CONGRUENCE HECKE OPERATORS » (title and author typed on p. 56;
%   headers « atkin page 2 », « atkin page 4 » on pp. 57 and 59, the header
%   cut off on pp. 58 and 60): §1 Introduction, §2 « The full modular
%   group », Hecke operators T_p, Theorem 1 (Hecke-Petersson), the
%   conjectured congruence analogue for c(n), Conjectures 1 and 2. Same
%   evidence (author's name, English, the author's own references to « (2) »
%   as his earlier paper). The copy continues into batch 4.
%   Publication, for the maintainers only and not checked here: the two
%   texts are believed to be Atkin's papers in Computers in Mathematical
%   Research (North-Holland, 1968) and Proc. Sympos. Pure Math. 12 (1969);
%   the notes name only author and title, which the leaves carry.
%
% The marks on these leaves: every handwritten mark is part of the
% photocopied image (same grey as the typing), not ink added on the copy,
% unlike the « (1968) » of p. 40. They are the author's hand-filled formulas
% and a copy-editor's English proof marks and instructions to the printer
% (« (comma) », « (hyphen) », « German 'pee' », « Cap. gamma », « lambda »,
% « alpha », « footnote on cover page », « online », « # »). None is
% attributable to Grothendieck, and the notes say so page by page, giving
% the worded ones. The one mark with a mathematical opinion, « multiplicat.
% semi group! » on p. 58 (against « the group of T_p »), is in the same
% hand as the printer's instructions (compare its « S » with « Sigma », p.
% 59); it is given in full, as a \note, not as \marginal, and not
% attributed to him.
%
% His pagination: none. Atkin's typed numbers are recorded per page; on
% pp. 45 and 46 the typed « 5 » and « 6 » are overwritten by hand « 6 » and
% « 7 » (the copy's own renumbering). No date on these leaves other than
% the years in Atkin's text (1962, 1963, 1964; references to 1967).
%
% What the batch is about: reference material only — Atkin's account of
% the computer search for congruences of the coefficients c(n) of j and of
% p(n) modulo powers of primes (Ramanujan-type versus multiplicative
% congruences, a « p-adic » Hecke theory), and the beginning of his
% « Congruence Hecke operators ». Nothing of his own.
%
% \ill{}: 0. \uncertain{}: 1 (the first word of the p. 58 margin note).

\documentclass[11pt,a4paper]{article}
% Le préambule commun à toutes les transcriptions du fonds Grothendieck.
%
% Il tient en une idée : **l'incertitude se compose, elle ne se résout pas.**
% Une transcription qui lisse un mot douteux a détruit la seule chose pour
% laquelle on la faisait — un lecteur doit pouvoir savoir, à chaque endroit,
% ce qui était sur le feuillet et ce qui a été ajouté. Les six macros qui
% suivent sont donc l'appareil critique, et non de la décoration.
%
% Les mêmes commandes sont reconnues par `scripts/render.mjs`, qui produit la
% vue de lecture du site. Une macro ajoutée ici doit l'être aussi là-bas,
% faute de quoi le rendu échouera bruyamment — ce qui est voulu : un rendu
% silencieusement faux est pire qu'un rendu absent.

\ProvidesPackage{grothendieck}[2026/08/08 Transcriptions du fonds Grothendieck]

\usepackage{amsmath,amssymb,amsthm}
\usepackage{mathrsfs}
\usepackage{stmaryrd}
% Les diagrammes commutatifs sont partout dans ce fonds : c'est la langue
% même de la géométrie algébrique de Grothendieck, et une transcription qui
% les rendrait en prose ne transcrirait rien.
\usepackage{tikz-cd}
% Français par défaut — c'est la langue des feuillets — mais l'anglais est
% chargé aussi : la traduction et le résumé sont des documents anglais, et la
% typographie française (espaces avant les deux-points) y serait une faute.
% Ces documents basculent par \selectlanguage{english} en tête de corps.
\usepackage[english,french]{babel}
\usepackage{fontspec}
\usepackage{geometry}
\geometry{a4paper,margin=25mm,marginparwidth=28mm}
\usepackage{marginnote}
\usepackage{xcolor}
% ulem, not soul. `soul`'s \st and \ul are built on a character-by-character
% scan that XeLaTeX's font machinery breaks: the document fails with an
% undefined control sequence rather than degrading. ulem does the same two jobs
% and is Unicode-engine safe. `normalem` stops it hijacking \emph, which the
% transcriptions use for Grothendieck's own emphasis.
\usepackage[normalem]{ulem}
\usepackage{hyperref}
\hypersetup{colorlinks=true,linkcolor=black,urlcolor=black,pdfborder={0 0 0}}

% --- Métadonnées du lot -----------------------------------------------------
% Reprises telles quelles par le script de rendu, qui les affiche en tête de
% la vue de lecture. Elles fixent ce que le fichier prétend couvrir : sans
% elles, une transcription flotte sans point d'attache dans un fonds de seize
% mille feuillets.

\newcommand{\folder}[1]{\gdef\grothFolder{#1}}
\newcommand{\batch}[1]{\gdef\grothBatch{#1}}
\newcommand{\leaves}[2]{\gdef\grothFirst{#1}\gdef\grothLast{#2}}
\newcommand{\foldertitle}[1]{\gdef\grothFoldertitle{#1}}
\gdef\grothFolder{?}\gdef\grothBatch{?}\gdef\grothFirst{?}\gdef\grothLast{?}\gdef\grothFoldertitle{}

% --- Le résumé --------------------------------------------------------------
%
% Il ouvre la lecture modernisée, sous le titre « Résumé », et il est composé
% en petit. Ce n'est pas un ornement : c'est le seul passage du document qui
% ne soit pas des mathématiques, et un lecteur doit pouvoir le reconnaître
% avant de l'avoir lu — puis le sauter s'il connaît déjà le sujet, ou s'y
% arrêter s'il ne le connaît pas. Le filet qui le suit marque la reprise du
% corps.

\newenvironment{resume}
  {\small\setlength{\parskip}{0.45em}}
  {\par\medskip\hrule\medskip}

% --- Mots-clés ---------------------------------------------------------------
%
% En anglais, en clôture du résumé de la lecture modernisée : le vocabulaire
% moderne sous lequel chercher ce que les feuillets construisent. Le manifeste
% du site (`npm run manifest`) les extrait pour étiqueter la cote entière —
% cette ligne est la source unique des tags, il n'y a pas de fichier à côté.
\newcommand{\keywords}[1]{\par\smallskip\noindent
  {\footnotesize\textsc{Keywords} --- \textit{#1}}\par}

% --- L'appareil critique ----------------------------------------------------

% Le numéro de feuillet, en marge. C'est l'ancre qui permet de régler un
% désaccord sur une formule en regardant *un* feuillet plutôt qu'une chemise
% de six cents. Le site s'en sert aussi pour tourner les pages du fac-similé
% quand on fait défiler la transcription.
\newcommand{\leaf}[1]{%
  \marginnote{\footnotesize\color{gray!70}\textsf{#1}}%
  \label{leaf:#1}\ignorespaces}

% Illisible : on ne devine pas. Un mot inventé qui se lit comme les autres est
% le pire résultat possible.
\newcommand{\ill}{\textcolor{red!60!black}{[\,\ldots\,]}}

% Lecture douteuse : le mot est donné, et signalé comme incertain.
\newcommand{\uncertain}[1]{\uline{#1}}

% Ajout de l'éditeur — une lettre manquante, un mot sous-entendu.
\newcommand{\add}[1]{\textcolor{blue!50!black}{[#1]}}

% Biffé par Grothendieck lui-même. Ce qu'il a rayé fait partie du document :
% une rature dit souvent où la pensée a changé de direction.
\newcommand{\struck}[1]{\sout{#1}}

% Note du transcripteur, sur l'état matériel du feuillet.
\newcommand{\note}[1]{\par\smallskip{\small\color{gray!60!black}\textit{[#1]}}\par\smallskip}

% Note marginale de Grothendieck, distincte du corps du texte.
\newcommand{\marginal}[1]{\par\smallskip\noindent\hspace{1em}%
  {\small\color{gray!60!black}#1}\par\smallskip}

% --- Titre ------------------------------------------------------------------

\AtBeginDocument{%
  \begin{center}
    {\large\bfseries Fonds Alexandre Grothendieck --- cote n\textsuperscript{o}~\grothFolder}\\[2pt]
    {\itshape\grothFoldertitle}\\[4pt]
    {\small feuillets \grothFirst--\grothLast\ (lot \grothBatch)}\\[2pt]
    {\footnotesize\color{gray!60!black}Transcription établie d'après le fac-similé numérisé
     par l'Université de Montpellier.\\
     Les lectures incertaines sont soulignées, les illisibles notées [\,\ldots\,],
     les ajouts entre crochets.}
  \end{center}
  \vspace{1em}}

\endinput


\folder{63}
\batch{3}
\pages{41}{60}
\dating{[à partir de 1964-vers 1970]}
\watermark{Édition de démonstration}
\foldertitle{Formulaire courbes elliptiques : notes et copies de notes manuscrites (s.d.), tapuscrit (s.d.), copies d'articles (s.d.)}

\begin{document}

\page{41}
\note{Les pages 41 à 54 sont la suite (pages dactylographiées « atkin 2 » à
« atkin 15 ») de la copie de l'article de A.~O.~L.~Atkin, « Congruences for
modular forms », commencée page 40. Texte d'un autre auteur : il n'est pas
recomposé ; chaque page reçoit une note. Les formules y sont complétées à la
main, d'une écriture anglaise ronde et droite qui n'est pas la sienne ; les
marques marginales sont des corrections d'épreuves et des instructions à
l'imprimeur, en anglais, qui font partie de l'image photocopiée. Aucune
marque de sa main n'est discernable sur ces pages.}
\note{Page « atkin 2 » : remerciements (Churchhouse, Fossey, Stokoe, système
Atlas Hartran) ; public visé ; §2 « Ramanujan congruences for primes
$q\leq 11$ » : la conjecture de Ramanujan (2.1), si
$24n-1\equiv 0 \pmod{5^a7^b11^c}$ alors $p(n)\equiv 0\pmod{5^a7^b11^c}$ ;
(2.2) $x=e^{2\pi i\tau}$, $f(x)=(1-x)(1-x^2)(1-x^3)\cdots$,
$\eta(\tau)=e^{\pi i\tau/12}f(x)$, $\sum p(n)x^n=1/f(x)$. Marques
d'épreuve : « (comma) » dans la marge droite, un trait oblique dans la
marge gauche.}

\page{42}
\note{Page « atkin 3 » : l'échec pour $7^3$ (Chowla, Gupta) ; Watson (1938),
Atkin (1967) ; forme correcte (2.3) avec $\beta=[(b+2)/2]$ ; Lehner et les
coefficients $c(n)$ de
$j(\tau)=\sum_{n\geq -1}c(n)x^n=\lbrace 1+240\sum\sigma_3(n)x^n\rbrace^3/xf^{24}(x)-744$
(2.4), $j$ Hauptmodul du groupe modulaire ; début de (2.5). Seules marques :
corrections d'épreuve (une lettre à insérer dans « propty », signalée dans
la marge gauche).}

\page{43}
\note{Page « atkin 4 » : fin de (2.5), dû à Lehner sauf les cas $e>3$
(Atkin 1967) ; §3 « Congruences for $q=13$ » : Newman (1958),
$c(13^2n)\equiv 8c(13n)\pmod{13}$ (3.1), et pour $t(n)=-c(13n)$,
$t(np)-t(n)t(p)+p^{11}t(n/p)\equiv 0\pmod{13}$ (3.2) ; (3.3)
$c(7\cdot 13^n)\equiv 0$, (3.4) $c(91n)\equiv 0\pmod{13}$ si $(n,7)=1$.
Aucune marque en dehors des formules manuscrites du texte.}

\page{44}
\note{Page « atkin 5 » : $t(n)\equiv\tau(n)\pmod{13}$,
$\sum\tau(n)x^n=xf^{24}(x)$, la relation (3.5)
$\tau(np)-\tau(n)\tau(p)+p^{11}\tau(n/p)=0$ (Ramanujan, Mordell), (3.6)
multiplicativité ; Glaisher (1907), Rankin. La page s'arrête à mi-hauteur.
Aucune marque en dehors du texte.}

\page{45}
\note{Page dactylographiée « atkin 5 », le 5 surchargé à la main en 6 :
§3.2, le résultat d'Atkin et O'Brien (3.7) :
$c(13^{\alpha+1}n)\equiv k_\alpha c(13^\alpha n)\pmod{13^\alpha}$ ; calculs
sur Atlas (1964) de $c(13^6n)$ modulo $13^6$ pour $n\leq 40960$, sous la
forme $a(n)\,13^{b(n)}$. Aucune marque en dehors du texte.}

\page{46}
\note{Page dactylographiée « atkin 6 », surchargée à la main en 7 : les
conjectures (3.8), puis (3.9)
$t(np)-t(n)t(p)+p^{-1}t(n/p)\equiv 0\pmod{13^\alpha}$ pour
$t(n)=c(13^\alpha n)/c(13^\alpha)$ ; (3.10) forme ramifiée
$t(13n)\equiv t(13)t(n)$ ; définition des « M-results » et
« R-results » ; (3.11) $c(13^\alpha P n)\equiv 0\pmod{13^\alpha}$ si
$(n,P)=1$. Marques d'épreuve dans la marge droite (une marque en face de (3.9),
« /n », « from (3.2) » à insérer après « show »).}

\page{47}
\note{Page « atkin 8 » : $P=\prod p_i^{\theta_i}$ ; divisibilité de $c(n)$
par $13^\alpha$ en densité positive ; le cas de $p(n)$ et des formes de
dimension demi-entière (13 « accidents ») ; §4 « The cases $q=17$, 19, and
23 » : les conjectures (3.9) et (3.10) valent mot pour mot pour ces $q$.
Marques d'épreuve dans la marge gauche (suppressions).}

\page{48}
\note{Page « atkin 9 » : l'opérateur $U_q$, $\sum c(q^\alpha n)x^n=U^\alpha
j(x)$ ; pour $q=19$, les fonctions $g_n(x)$ et le système (4.1)
$U_{19}\,j\equiv a_1g_1+19a_2g_2+19a_3g_3\pmod{19^2}$, etc. ; certaines
constantes $c_{ij}$ contiennent des puissances de 19 non déductibles
\emph{a priori}. Marques d'épreuve : « etc. » biffé, « $a_i$ and » ajouté,
« (hyphen) ».}

\page{49}
\note{Page « atkin 10 » : les programmes s'auto-vérifient ; genres $0,1,1,2$
de $\Gamma_0(13)$, $\Gamma_0(17)$, $\Gamma_0(19)$, $\Gamma_0(23)$ ; §5 « The
cases $q\geq 29$ » : le système modulo 29 et la congruence
$\mu^2-8\mu+5\equiv 0\pmod{29}$ sans racine, d'où le passage à
$\mathrm{GF}(29^2)$. Marques d'épreuve : « (hyphen) », « not » à insérer.}

\page{50}
\note{Page « atkin 11 » : $c(29n)=a\lambda(n)+\bar a\bar\lambda(n)$
dans $\mathrm{GF}(29^2)$ et (5.2), (5.3), la conjecture générale (5.4) pour
les puissances de 29, prouvée pour $\alpha\leq 5$ ; $q=31$ analogue,
$q=37$ demande une extension cubique ; pour $41\leq q\leq 67$, $c(q^\alpha
n)$ combinaison de $[q/12]$ coefficients multiplicatifs. Marques d'épreuve :
« (GF) » entouré, « multiplicative » inséré, « l/ ».}

\page{51}
\note{Page « atkin 12 » : §6 « The analogy with Hecke theory » : l'exemple
de Hecke $\sum t_2(n)x^n=x^2f^{48}(x)$, $t_2(n)=a\lambda(n)+\bar a\bar
\lambda(n)$ dans $\mathbf{Q}(\sqrt{144169})$, (6.1) ; conjecture d'une
théorie de Hecke « $\mathfrak{p}$-adique » pour les formes de dimension
positive ou nulle ; §7 « Non-Ramanujan congruences for primes $q\leq 11$ ».
Marques d'épreuve : « (German p) », « $\tau=$ » à insérer deux fois devant
« $i\infty$ », lettres à insérer.}

\page{52}
\note{Page « atkin 13 » : exemple $q=5$, $t(n)=c(5^\alpha n)/c(5^\alpha)$,
(7.1) ; ramanujaniennes et multiplicatives ; §8 « Notes on computation » :
coût $O(N)$ des opérations sur les séries. Marques d'épreuve : « online »
en face de la note de bas de page, lettres et « $x^n$ » à insérer.}

\page{53}
\note{Page « atkin 14 » : $O(\sum_{p\leq N}N/p)=O(N\log\log N)$ ;
multiplication par la série d'Euler $f(x)=\prod(1-x^r)$, en $O(N^{3/2})$ ;
limites de taille ; série de Jacobi $f^3(x)$ ; la mémoire comme facteur
limitant sur Atlas. Une marque d'épreuve dans la marge gauche.}

\page{54}
\note{Page « atkin 15 », dernière de l'article dans cette copie : mémoire,
programmes Fortran, la fonction IJMODK de Don Russell, et une panne
matérielle révélée en 1964. Marques d'épreuve (« e/ » deux fois, « timing »
inséré). La page s'arrête au premier tiers ; pas de bibliographie.}

\page{56}
\note{Les pages 56 à 60 sont les cinq premiers feuillets d'une copie
dactylographiée d'un second texte de A.~O.~L.~Atkin, « Congruence Hecke
Operators » (titre et auteur dactylographiés en tête de cette page). Texte
d'un autre auteur : non recomposé, une note par page. Mêmes caractères que
pour l'article précédent : formules complétées à la main par l'auteur,
corrections d'épreuve et instructions à l'imprimeur en anglais, toutes dans
l'image photocopiée ; aucune marque qui soit de sa main.}
\note{Page 1 : §1 « Introduction » : l'article (2) de l'auteur a suivi
l'ordre de la découverte ; celui-ci part de l'autre bout et esquisse une
théorie « $\mathfrak{p}$-adique » des opérateurs de Hecke ; les
congruences de Ramanujan pour $p(5n+4)$, $p(7n+5)$, $p(11n+6)$ et
$p(206839n+2623)\equiv 0\pmod{17}$. Instructions à l'imprimeur : « German
'pee' », « Cap. gamma ».}

\page{57}
\note{Page « atkin page 2 » : §2 « The full modular group » : $\Gamma$,
$x=e^{2\pi iz}$, formes modulaires de dimension $k$, conditions (M1), (M2),
(M3a-c) ; formes entières, régulières, paraboliques, espaces $C$, $C^+$,
$C^0$ ; série de Fourier $F(x)=\sum_{n\geq N}a(n)x^n$. Instructions :
« pi », « (are) », suppressions.}

\page{58}
\note{Page dactylographiée 3 (en-tête coupé) : l'opérateur de Hecke
$T_pF(x)=\sum\lbrace a(np)+p^{-k-1}a(n/p)\rbrace x^n$, propriétés (H1),
(H2), produit de Petersson, (H3), (H4) ; Théorème 1 (Hecke-Petersson) :
base de formes propres, relation (1)
$a(np)-a(n)a(p)+p^{-k-1}a(n/p)=0$. Instructions à l'imprimeur : « lambda »,
« alpha », « xi eta », « footnote on cover page », un signe dièse (espace). Dans la marge
gauche, en face de « the group of $T_p$ » (« group » biffé), trois mots
d'une écriture rapide qui paraît être celle des instructions à
l'imprimeur (comparer le « S » de « Sigma », page 59), et non la sienne :
« \uncertain{multiplicat.}\ semi group! ».}

\page{59}
\note{Page « atkin page 4 » : l'extension du Théorème 1 à $C$ échoue par
(H4) ; le choix de $j(z)=x^{-1}+744+\sum c(n)x^n$ et la question d'une
propriété (2) $a(1)=1$, $a(np)-a(n)a(p)+p^{-1}a(n/p)=0$ pour des $a(n)$
dans $\mathbf{Z}/q^\alpha\mathbf{Z}$. Instructions : « Sigma », « alpha ».}

\page{60}
\note{Page dactylographiée 5 (en-tête coupé) : Conjecture 1 ($q\leq 23$,
$a(n)=c(q^\alpha n)/c(q^\alpha)$ multiplicatif modulo $q^\alpha$) ;
congruences (3) de type Ramanujan ; Conjecture 2 ($q=29$ ou 31, extension
quadratique $E$ de $\mathbf{Z}/q^\alpha\mathbf{Z}$) ; le cas $q\geq 37$ ;
passage à $\Gamma_0(q)$. Instructions : « theta », « beta », un signe dièse (espace),
mots à insérer. La copie se poursuit dans le lot suivant.}

\end{document}
